\section{What we do}


From hereon, we rely on a few key assumptions:

\begin{enumerate}
\item We assume that the person studying the code has \emph{limited knowledge} of the code's internals, the application domain, and the algorithms used.
\item We evaluate \emph{how} a code performs and whether certain characteristics deserve closer examination.
\item We are interested in the code's behavior on \emph{one particular machine} for \emph{one particular experimental setup} or a closely defined set of experiments.
\item Whenever we explain \emph{why} a code performs as it does, this
explanation comes from a machine's point of view: It does not argue about the
algorithm or science case.
\end{enumerate}

\begin{readerguide}
 Feel free to skip the remainder of this chapter if you are a first-time reader. 
\end{readerguide}


\noindent
In scientific computing, we typically work with an iterative approach:
First, we \emph{measure} or \emph{benchmark} how a code performs.
Second, we use these data to inform our \emph{performance assessment}, i.e.~our reporting.
Third, we conduct a \emph{performance analysis} to explain why the code exhibits the observed behaviour.
Finally, this analysis feeds into program tuning or optimisation, and we repeat the workflow.



If we take our list of activities and integrate it within the definition of the performance analysis workflow, it becomes clear that the tuning step is out-of-scope here.
Another related technique is out-of-scope, too:
The most advanced papers in scientific computing, particularly within high-performance computing, deliver a performance model;
essentially a pen-and-paper assessment of what a code should be able to deliver.
This is impossible if we adhere to the principles above: having limited
knowledge and only explaining observed behavior rather than relating it to
external factors (such as input data or algorithmic choices) or algorithmic
context.
We commit to a data-driven approach.
Obviously, such an approach might complement an analytic, theoretical assessment
of code, but we recommend considering this as an entirely separate cup of tea.



We will inevitably face difficulties if we truly do not know the code we are evaluating:
Explaining certain effects is so much easier with some (rudimentary)
understanding of the code.
However, there are three compelling reasons to treat code as a (logical) black box:
First, performance assessment is often delegated to non-developers: colleagues
from computing centers, centralised HPC specialists, or external service providers.
If we want to discuss performance analysis methodology, it is counterproductive to assume in-depth code knowledge.
Many people conducting assessments cannot be experts in every code they examine.
Second, requiring code expertise within assessment teams would imply that core developers are always available.
This might be an ideal in the spirit of eXtreme Programming \cite{beck1999}, where
the customer must be accessible at all times.
It is unrealistic in science, where PIs are busy with various tasks, developers
frequently leave projects---as they graduate, accept new positions, or are
reassigned---and where code components are often written by individuals and are
so complex that only a few people understand the core parts.
Finally (and most importantly), knowledge about a codebase tends to bias the assessment.




Developers inevitably make assumptions about reasons for code behaviour when working on performance analysis.
(Senior) Developers tend to ``defend'' their code, explaining any flaws as
natural consequences of their domain's challenging nature, their awareness of future requirements, or their in-depth knowledge of other setups not covered by the present configuration.
Simultaneously, we have frequently observed junior developers downplaying or
challenging existing code---likely to prove themselves by pointing out ``how
things really should work''---and consequently interpreting every piece of data
through this lens.
This list is not comprehensive.
The whole team dynamics dimension deserves further discussion (in
Section~\ref{section:communication:team-dynamics}).
For the time being, we conclude that having emotionally invested developers
involved in early assessment is counterproductive due to their inherent bias.


\begin{assessmentcrime}
\label{assessmentcrime:code-expert-bias}
People conduct assessments as code or performance experts.
Due to this, they deliver biased assessments.
\end{assessmentcrime}


For a high-quality assessment, it helps to step back and (pretend to) know nothing.
This way, we avoid entering performance analysis with bias.
Adhere strictly to the observational perspective.
Even if it is your own code, try to detach yourself and pretend you have no prior knowledge.
Abstract.





\paragraph*{How it works}

\begin{itemize}[noitemsep]
  \item[$\boxplus$] Pretend you know nothing about the code. Blend out your
  expert knowledge and start on a fresh sheet.
  \item[$\boxplus$] Stick to the data you measure. 
  \item[$\boxplus$] Ignore what others pretend to know or have found. 
\end{itemize}


\paragraph*{How it does not work}

\begin{itemize}[noitemsep]
  \item[$\boxminus$] Start with the attitude ``I know what this code does and
  where there are weaknesses''.
  \item[$\boxminus$] Skip assessment steps, as you already know what's going on.
  \item[$\boxminus$] Involve someone who's emotionally involved in the code's
  development into the reporting. They will always defend their brain child.
\end{itemize}

